\chapter{注解与反射}
\label{ch:20}

第 19 章的预处理改的是源码文本。这一章讲一个更「上层」的能力：注解——在声明旁边贴一张标签，标签上带点数据。平时它安安静静躺着不干活；但程序可以在运行期把标签读回来，据此决定怎么行动。

一句话概括：注解 = 给代码贴的、可以被程序自己读到的元数据。

\section{声明一个注解类型}

注解类型用 \texttt{@annotation} 声明，参数写成「名字: 类型」，可以带默认值：

\par
\medskip
\begin{lstlisting}[style=chali]
@annotation Route(path: String, level: int32 = 2)
\end{lstlisting}
\par
\medskip

读作：声明一个叫 \texttt{Route} 的注解，它有两个参数 \texttt{path}（字符串）和 \texttt{level}（整数，缺省 2）。

\section{把注解贴到声明上}

声明完，就能贴了——写在被标注的东西前面（下面这段要在反射版工具链下跑，见 20.6）：

\par
\medskip
\begin{lstlisting}[style=chali]
@annotation Route(path: String, level: int32 = 2)

@Route("/health", level: 1)
int64 health()
{
    return 200;
}

int64 main()
{
    println("{}", health());
    return 0;
}
\end{lstlisting}
\par
\medskip

\par
\medskip
\begin{lstlisting}[style=out]
200
\end{lstlisting}
\par
\medskip

注意：这个程序的行为一点没变——\texttt{health(\allowbreak{})\allowbreak{}} 还是返回 200。注解本身不影响执行，它只是一份「挂上去的资料」。要让它起作用，得靠 20.4 节之后的手段去读它。

参数可以按位置写（\texttt{"/\allowbreak{}health"}），也可以按名字写（\texttt{level:\allowbreak{} 1}）。

\section{能贴在哪里}

注解几乎能贴在所有声明上：

\begin{center}
\begin{tabular}{ll}
\toprule
位置 & 例子 \\
\midrule
函数 & \texttt{@Route(\allowbreak{}"/\allowbreak{}health")\allowbreak{} int64 health(\allowbreak{})\allowbreak{} \{ … \}} \\
结构体 / 接口 / 枚举 & \texttt{@Component struct Service \{ … \};} \\
字段 & \texttt{@Note(\allowbreak{}"field")\allowbreak{} int64 value;} \\
方法 / 构造器 & \texttt{@Route(\allowbreak{}"/\allowbreak{}read")\allowbreak{} int64 read(\allowbreak{})\allowbreak{} \{ … \}} \\
静态成员 & \texttt{@Note(\allowbreak{}"static")\allowbreak{} static int64 shared;} \\
形参 & \texttt{@Note(\allowbreak{}"contract")\allowbreak{} int64 read(\allowbreak{})\allowbreak{};}（接口方法声明） \\
\bottomrule
\end{tabular}
\end{center}

多个注解可以叠着贴：

\par
\medskip
\begin{lstlisting}[style=chali]
@Component
@Route("/read", level: 3)
int64 read()
{
    return self.value;
}
\end{lstlisting}
\par
\medskip

\section{读注解：运行期反射}

注解真正的用处，是程序在运行期把标签读回来。这套能力叫反射，在标准模块 \texttt{std:\allowbreak{}:\allowbreak{}reflect} 里，要 \texttt{import}：

\par
\medskip
\begin{lstlisting}[style=chali]
import std::reflect;

@Reflect
struct Point
{
    int64 x;
    int64 y;
};

int64 main()
{
    Option<std::reflect::StructInfo> info = std::reflect::getStructInfo<Point>();
    println("{}", info.isValue());
    std::reflect::StructInfo s = info.unwrap();
    println("{} {}", s.name, s.fields.len());
    return 0;
}
\end{lstlisting}
\par
\medskip

\par
\medskip
\begin{lstlisting}[style=out]
1
Point 2
\end{lstlisting}
\par
\medskip

几个要点：

\begin{itemize}

\item \texttt{@Reflect} 是「暴露」标记：只有被 \texttt{@Reflect} 标过的声明，运行期才查得到。没标的类型，\texttt{get\allowbreak{}Struct\allowbreak{}Info} 会返回 \texttt{None}（不是报错）；

\item \texttt{get\allowbreak{}Struct\allowbreak{}Info<\allowbreak{}Point>\allowbreak{}(\allowbreak{})\allowbreak{}} 用类型参数取元数据，返回 \texttt{Option<\allowbreak{}Struct\allowbreak{}Info>\allowbreak{}}（第 16 章）；

\item \texttt{Struct\allowbreak{}Info} 里写着 \texttt{name}、\texttt{fields}、\texttt{methods}、\texttt{static\allowbreak{}Members}、\texttt{annotations}……——一个类型能问的都在里头。

\end{itemize}

读注解的实参：

\par
\medskip
\begin{lstlisting}[style=chali]
import std::reflect;

@annotation Route(path: String, level: int32 = 2)

@Reflect
@Route(path: "/health")
int64 health()
{
    return 200;
}

int64 main()
{
    List<std::reflect::Annotation> anns = std::reflect::getFuncAnnotations("health");
    println("{} {}", anns.len(), anns[1].args.len());
    println("{}", anns[1].args[0].asText().unwrap());
    return 0;
}
\end{lstlisting}
\par
\medskip

\par
\medskip
\begin{lstlisting}[style=out]
2 2
/health
\end{lstlisting}
\par
\medskip

\begin{itemize}

\item \texttt{get\allowbreak{}Func\allowbreak{}Annotations(\allowbreak{}"函数名")\allowbreak{}} 按名字取函数上的注解列表；

\item 这一份列表里有两项：\texttt{@Reflect}（0 个实参）和 \texttt{@Route}（2 个实参）；

\item 实参是一串 \texttt{std:\allowbreak{}:\allowbreak{}reflect:\allowbreak{}:\allowbreak{}Reflect\allowbreak{}Value}（一个能装任意类型的小盒子），用 \texttt{as\allowbreak{}Text(\allowbreak{})\allowbreak{}} / \texttt{as\allowbreak{}I32(\allowbreak{})\allowbreak{}} 等按类型取出来，都返回 \texttt{Result}。

\end{itemize}

\begin{tip}{提示}
注解实参会补全默认值：\texttt{Route} 的 \texttt{level} 缺省是 2，所以 \texttt{args} 里 \texttt{level} 位永远有值。
\end{tip}

\section{按名字调用函数}

反射还能按字符串名字调函数——名字在运行期才知道时，只有这条路：

\par
\medskip
\begin{lstlisting}[style=chali]
import std::reflect;

@Reflect
int64 plusOne(int64 x)
{
    return x + 1;
}

int64 main()
{
    List<std::reflect::ReflectValue> args;
    args.push(std::reflect::ReflectValue::I64(41));
    Result<std::reflect::ReflectValue, std::reflect::ReflectError> r = std::reflect::call("plusOne", args);
    println("{}", r.isOk());
    return 0;
}
\end{lstlisting}
\par
\medskip

\par
\medskip
\begin{lstlisting}[style=out]
1
\end{lstlisting}
\par
\medskip

实参要自己手工装箱（\texttt{Reflect\allowbreak{}Value:\allowbreak{}:\allowbreak{}I64(\allowbreak{}41)\allowbreak{}}），返回值也是 \texttt{Result<\allowbreak{}Reflect\allowbreak{}Value,\allowbreak{} …>\allowbreak{}}。

调方法要用 \texttt{call\allowbreak{}Method<\allowbreak{}T>\allowbreak{}}，并且把接收者写成引用：

\par
\medskip
\begin{lstlisting}[style=chali]
import std::reflect;

@Reflect
struct Counter
{
    int64 n;
    int64 bump()
    {
        self.n = self.n + 1;
        return self.n;
    }
};

int64 main()
{
    Counter c;
    c.n = 0;
    List<std::reflect::ReflectValue> args;
    Result<std::reflect::ReflectValue, std::reflect::ReflectError> r =
    std::reflect::callMethod<Counter>(&c, "bump", args);
    println("{} {}", r.isOk(), c.n);
    return 0;
}
\end{lstlisting}
\par
\medskip

\par
\medskip
\begin{lstlisting}[style=out]
1 1
\end{lstlisting}
\par
\medskip

按名的接口分派有 \texttt{call\allowbreak{}Interface<\allowbreak{}I>\allowbreak{}}——写法一样，把接口名当类型参数。

\section{构建与目标 VM 的反射能力}

根目录工具链默认同时编入前端和运行时反射，单次编译也默认启用。本例使用默认工具链即可运行：

\par
\medskip
\begin{lstlisting}[style=chali]
@annotation Route(path: String, level: int32 = 2)

@Route("/health", level: 1)
int64 health()
{
    return 200;
}

int64 main()
{
    println("{}", health());
    return 0;
}
\end{lstlisting}
\par
\medskip

\par
\medskip
\begin{lstlisting}[style=out]
200
\end{lstlisting}
\par
\medskip

{\small\setlength{\tabcolsep}{4pt}
\begin{longtable}{>{\raggedright\arraybackslash}p{\dimexpr 0.151\linewidth-2\tabcolsep\relax}>{\raggedright\arraybackslash}p{\dimexpr 0.829\linewidth-2\tabcolsep\relax}}
\toprule
方式 & 含义 \\
\midrule
\endhead
\texttt{make} & 根目录把前端 CHALI\_WITH\_REFLECTION=1、运行时 STACKVM\_WITH\_REFLECTION=1 钉死：反射不可裁剪 \\
\texttt{make reflection} & 同义的显式入口（两项本来就固定为 1） \\
不写单次开关 & 默认写入反射元数据 \\
\texttt{--reflection} & 显式要求写入反射元数据（缺省即如此） \\
\texttt{--no-reflection} & 不写反射元数据；关闭后使用注解或反射会报错。它只影响这次产物，不减少编译期工作量 \\
\bottomrule
\end{longtable}
}

\par
\medskip
\begin{lstlisting}[style=sh]
chali run app.chl
chali svmb app.chl --no-reflection --out build/board-app
\end{lstlisting}
\par
\medskip

第二条用于无需反射的源程序：产物必须匹配目标 VM 的反射构建配置。开启反射的产物携带特征位，关闭能力的 VM 会拒绝加载。若源程序依赖反射，应为目标编入能力并评估元数据和运行时预算，不能通过省略 \texttt{--reflection} 假装关闭。独立生成的编译器工程按其 CMake 配置单独启用两项宏。

\section{什么时候该用注解}

注解与反射提供运行期元数据和受限动态调用，会增加产物、分配和调用开销。默认桌面工具链支持它，裁剪目标需要显式评估能力和预算。

{\small\setlength{\tabcolsep}{4pt}
\begin{longtable}{>{\raggedright\arraybackslash}p{\dimexpr 0.384\linewidth-2\tabcolsep\relax}>{\raggedright\arraybackslash}p{\dimexpr 0.596\linewidth-2\tabcolsep\relax}}
\toprule
你的想法 & 该不该用 \\
\midrule
\endhead
「想给函数打点标记，运行时按标记分派」 & \textcolor{green!50!black}{\ding{51}} 典型场景（路由、插件注册） \\
「想让代码生成器读元数据」 & \textcolor{green!50!black}{\ding{51}} 合理 \\
「想根据标记做一遍扫描」 & \textcolor{green!50!black}{\ding{51}} 合理 \\
「想省一个函数参数」 & \textcolor{red!80!black}{\ding{55}} 用参数 \\
「想实现多态」 & \textcolor{red!80!black}{\ding{55}} 用接口 / 泛型（第 13、14 章） \\
「目标 VM 关闭反射」 & 移除反射依赖并用 --no-reflection 生成（反射本身仍不可裁剪） \\
\bottomrule
\end{longtable}
}

经验法则：注解是「框架和工具用的」，不是「业务逻辑用的」。写业务代码时，几乎总有更直接的办法。

\section{小结}

\begin{center}
\begin{tabular}{ll}
\toprule
你会了 & 关键点 \\
\midrule
声明注解 & \texttt{@annotation 名字(\allowbreak{}参数:\allowbreak{} 类型 =\allowbreak{} 默认值)\allowbreak{}} \\
贴注解 & 写在声明前；可叠多个；能贴允许的函数/类型/字段/方法/静态成员；不能贴形参 \\
注解不改变行为 & 它只是一份挂在声明上的资料 \\
\texttt{@Reflect} & 暴露标记：没标过的声明，运行期查不到 \\
查类型 & \texttt{std:\allowbreak{}:\allowbreak{}reflect:\allowbreak{}:\allowbreak{}get\allowbreak{}Struct\allowbreak{}Info<\allowbreak{}T>\allowbreak{}(\allowbreak{})\allowbreak{}} \ensuremath{\rightarrow} \texttt{Option<\allowbreak{}Struct\allowbreak{}Info>\allowbreak{}} \\
读注解 & \texttt{get\allowbreak{}Func\allowbreak{}Annotations(\allowbreak{}"名字")\allowbreak{}} / \texttt{get\allowbreak{}Annotations<\allowbreak{}T>\allowbreak{}(\allowbreak{})\allowbreak{}} \\
取实参 & \texttt{Annotation.\allowbreak{}args} 是 \texttt{Reflect\allowbreak{}Value}，用 \texttt{as\allowbreak{}Text(\allowbreak{})\allowbreak{}} 等取（返回 \texttt{Result}） \\
按名调用 & \texttt{std:\allowbreak{}:\allowbreak{}reflect:\allowbreak{}:\allowbreak{}call} / \texttt{call\allowbreak{}Method<\allowbreak{}T>\allowbreak{}(\allowbreak{}\&obj,\allowbreak{} …)\allowbreak{}} / \texttt{call\allowbreak{}Interface<\allowbreak{}I>\allowbreak{}} \\
构建 & 反射为固定组成、不可裁剪；无反射目标显式 --no-reflection（只影响产物） \\
取舍 & 框架/工具才用；业务逻辑优先接口/泛型 \\
\bottomrule
\end{tabular}
\end{center}

\section{下一步}

到这里，语言本身讲完了——从变量到注解，你能写出完整的、有类型安全、有资源管理、有抽象能力的 ChaLi 程序。

第七部分开始「走出去」：让 ChaLi 跟外部世界打交道。第一个是与 C 互操作——调用现成的 C 库、把自己的函数交给 C 调用。

翻到\hyperref[ch:21]{第 21 章 与 C 互操作}。

\section{元数据之后：精确类型的字段读写}

查询到字段名，并不意味着可以把它按任意类型取出。\texttt{get\allowbreak{}Field<\allowbreak{}T,\allowbreak{} F>\allowbreak{}} 和 \texttt{set\allowbreak{}Field<\allowbreak{}T,\allowbreak{} F>\allowbreak{}} 同时指定模型与字段类型；字段不存在、类型不同或模型未暴露时检查错误分支。

\par
\medskip
\begin{lstlisting}[style=chali]
import std::reflect;

@Reflect
struct Profile
{
    String name;
    int64 score;
};

int64 main()
{
    Profile profile;
    profile.name = "原名";
    profile.score = 80;
    assert(std::reflect::setField<Profile, String>(&profile, "name", "新名").unwrap());
    assert(std::reflect::getField<Profile, String>(&profile, "score").isErr());
    assert(std::reflect::getField<Profile, int64>(&profile, "missing").isErr());
    int64 score = std::reflect::getField<Profile, int64>(&profile, "score").unwrap();
    println("{} {}", profile.name, score);
    return 0;
}
\end{lstlisting}
\par
\medskip

\par
\medskip
\begin{lstlisting}[style=out]
新名 80
\end{lstlisting}
\par
\medskip

读取得到字段的独立值，修改模型通过写入接口进行。对容器、嵌套值和资源的支持边界仍须遵守实际接口；反射不会将不可复制资源自动转换成普通可复制值。

本章负责说明元数据及访问规则；模型与 JSON/TOML、字节序列化和动态 SQL 的组合见\href{https://gitee.com/rabbitc999/chali/blob/master/doc/36-reflection-serialization-sql.md}{反射编解码与动态 SQL}。

\section{边界与验证}

注解提供元数据，不自动实现业务逻辑。模型未暴露、字段缺失或类型不匹配，与编译器裁剪掉反射能力是不同情况；高级入口的 panic 边界按实际构建约定处理。

\section{习题}

前两题用于理解和基本实现，第三题检查边界，最后一题为综合挑战。先独立完成，再核对参考解答。

\begin{enumerate}

\item 加了注解是否就等于实现了注解名所描述的业务？

\item 定义 @Reflect 模型，按字段名读取一个整数。

\item 分别测试不存在字段和错误字段类型。

\item 挑战：对同一模型执行 JSON 与字节往返，检查副本独立性。

\end{enumerate}

\subsection*{参考解答与验收要点}

\begin{enumerate}

\item 不等于；必须有读取该元数据并执行相应规则的代码。

\item 使用精确字段类型并检查 Result，验证返回值等于原字段。

\item 两者都应进入错误路径，且不会修改模型或进行隐式转换。

\item 比较关键值，再修改解码结果；源对象必须保持原值，错误载荷应明确拒绝。

\end{enumerate}

\section{番外：榜上有名，方可查验}

改编自《封神演义》封神榜，以下为编程主题的独立改写。

姜子牙开榜，众人排队查询各自的职司。守榜人照登记册按名取值，办得很快。忽然有人指着榜外一名过客，说他既然也长着手脚，应当能查到同样字段。

守榜人要看暴露的元数据，那人却递来一张画像。画像再像，也没有字段名称、类型与注解。另一个人找到了字段，仍执意把数字按字符串读取，声称自己反正会认字。

姜子牙将榜、画像和实际任职分开放好：榜告诉你允许查什么，类型告诉你如何理解取出的值，任职是否合理还须另作判断。仅有一条登记记录，不会自动完成业务校验。

临近散场，守榜人发现整座查验台根本没有装进这次裁剪后的工具。姜子牙便让入口直接报告能力缺失，不能再拿“此人不在榜上”混作同一种错误。队伍这才知道，该补的是模型标记，还是整个工具的能力。

\subsection*{技术对照}

注解提供元数据，不自动实现业务逻辑。模型未暴露、字段缺失或类型不匹配，与编译器裁剪掉反射能力是不同情况；高级入口的 panic 边界按实际构建约定处理。
